%=====================================================================
\chapter{Relational Algebra}\label{ch:algebra}
%=====================================================================
\locator{4}{Part IV --- Asking Questions}

\begin{outcomes}
By the end of this chapter you will be able to:
\begin{tight}
  \item \textbf{Apply} selection, projection, Cartesian product, union,
        difference and rename to small relations by hand.
  \item \textbf{Write} an algebra expression for a question stated in English.
  \item \textbf{Translate} between an algebra expression and the
        \sqlkw{SELECT} statement that computes it.
  \item \textbf{Explain} why the output of every operator is a relation, and
        what that makes possible.
  \item \textbf{State} which operators need union-compatible inputs, and
        \textbf{check} that two relations are.
\end{tight}
\end{outcomes}

\begin{prereq}
\chref{relational}: a relation is a set of tuples. Appendix~B for set union,
intersection and difference --- if those three words are not routine, read it
first; it is four pages. This chapter is the first in the course with a
symbol on most pages.
\end{prereq}

\begin{keyterms}
relational algebra $\bullet$ operator $\bullet$ closure $\bullet$ selection
($\sigma$) $\bullet$ projection ($\pi$) $\bullet$ Cartesian product
($\times$) $\bullet$ union ($\cup$) $\bullet$ set difference ($-$) $\bullet$
intersection ($\cap$) $\bullet$ rename ($\rho$) $\bullet$ union-compatible
$\bullet$ predicate $\bullet$ query tree $\bullet$ procedural $\bullet$
declarative
\end{keyterms}

\section{Why a Course on SQL Starts With Algebra}

SQL is the language you will use. Relational algebra is the language the
\emph{database} uses --- the optimiser of \chref{whodoeswhat} turns your
\sqlkw{SELECT} into an algebra expression, rearranges it, and runs that.
Knowing the algebra is knowing what your query becomes.

It also answers a question that SQL hides: why can a query be used inside
another query? The answer is one property, and it is the first definition.

\begin{definitionbox}[Relational algebra, and closure]
The \term{relational algebra} is a set of operators. Each takes one or two
relations as input and produces \textbf{a relation} as output.

\smallskip
That last fact is \term{closure}: because every result is a relation, every
result can be the input to another operator. Expressions nest without limit,
exactly as $(3+4)\times 2$ does in arithmetic --- and for the same reason, which
is that $3+4$ is a number.

\smallskip
Six operators are fundamental: $\sigma$, $\pi$, $\times$, $\cup$, $-$ and
$\rho$. Every other operator, including every kind of join, can be written in
terms of these six.
\end{definitionbox}

\section{The Two Unary Operators}

\begin{definitionbox}[Selection and projection]
\term{Selection} $\sigma_{p}(R)$ returns the tuples of $R$ for which the
predicate $p$ is true. It chooses \textbf{rows}. The result has the same
attributes as $R$ and at most as many tuples.

\smallskip
\term{Projection} $\pi_{A_1,\dots,A_k}(R)$ returns the tuples of $R$
restricted to the attributes $A_1,\dots,A_k$. It chooses \textbf{columns}.
Because the result is a set, \textbf{duplicates are removed} --- so it may have
fewer tuples than $R$.
\end{definitionbox}

\dsfig{%
\begin{tikzpicture}[
  cell/.style={draw=primary!35, line width=0.5pt, minimum width=1.55cm,
               minimum height=0.42cm, font=\scriptsize\ttfamily,
               text=primary!65!black, inner sep=0pt},
  hd/.style={cell, fill=primary, text=white, font=\scriptsize\bfseries},
  sel/.style={cell, fill=greenhl!18},
  prj/.style={cell, fill=secondary!18},
  ttl/.style={font=\scriptsize\bfseries, text=#1, anchor=south},
  nt/.style={font=\scriptsize\itshape, text=neutral, align=center}
]
  % --- selection: rows
  \node[ttl=greenhl!60!black] at (-3.9,0.55)
       {$\sigma_{\,\mathit{programme}=\mathrm{'BSIT'}}(\mathit{student})$};
  \foreach \c/\h [count=\i] in {roll\_no/r, name/n, programme/p}{
    \node[hd] at (-5.45+\i*1.55-1.55,0.21) {\c};}
  \foreach \r/\a/\b/\c/\s in {1/IT-001/Ayesha/BSIT/sel, 2/IT-002/Bilal/BSIT/sel,
                              3/CS-007/Hira/BSCS/cell, 4/IT-014/Saad/BSIT/sel}{
    \node[\s] at (-5.45,-0.42*\r+0.21) {\a};
    \node[\s] at (-3.90,-0.42*\r+0.21) {\b};
    \node[\s] at (-2.35,-0.42*\r+0.21) {\c};}
  \node[nt, text width=4.6cm] at (-3.9,-2.25)
       {chooses \textbf{rows}; every column survives};

  % --- projection: columns
  \node[ttl=secondary!70!black] at (3.1,0.55)
       {$\pi_{\,\mathit{programme}}(\mathit{student})$};
  \foreach \c [count=\i] in {roll\_no, name, programme}{
    \node[hd] at (1.55+\i*1.55-1.55,0.21) {\c};}
  \foreach \r/\a/\b/\c in {1/IT-001/Ayesha/BSIT, 2/IT-002/Bilal/BSIT,
                           3/CS-007/Hira/BSCS, 4/IT-014/Saad/BSIT}{
    \node[cell] at (1.55,-0.42*\r+0.21) {\a};
    \node[cell] at (3.10,-0.42*\r+0.21) {\b};
    \node[prj]  at (4.65,-0.42*\r+0.21) {\c};}
  \draw[-{Stealth[length=2mm]}, secondary!70, line width=0.9pt]
       (5.55,-0.6) -- (6.35,-0.6);
  \node[hd]   at (7.2,-0.21) {programme};
  \node[prj]  at (7.2,-0.63) {BSIT};
  \node[prj]  at (7.2,-1.05) {BSCS};
  \node[nt, text width=4.8cm] at (3.1,-2.25)
       {chooses \textbf{columns}, then removes duplicates: four rows in, two
        out};
\end{tikzpicture}}%
{The two unary operators. Selection cuts horizontally and never loses a
column; projection cuts vertically and can lose rows, because three
\texttt{BSIT} values become one.}{selproj}

\begin{examplebox}[Both operators, and their SQL]
\begin{tabular}{@{}L{5.6cm}L{8.4cm}@{}}
\textbf{Algebra} & \textbf{SQL} \\
$\sigma_{\mathit{semester}=2}(\mathit{student})$ &
  \texttt{SELECT * FROM student WHERE semester = 2;} \\
$\pi_{\mathit{name},\,\mathit{programme}}(\mathit{student})$ &
  \texttt{SELECT DISTINCT name, programme FROM student;} \\
$\pi_{\mathit{name}}(\sigma_{\mathit{semester}=2}(\mathit{student}))$ &
  \texttt{SELECT DISTINCT name FROM student WHERE semester = 2;} \\
\end{tabular}

\smallskip
\textbf{Look at the word \sqlkw{DISTINCT}.} It is in all three projection rows
because algebra projection removes duplicates and SQL's does not. SQL returns
a \emph{bag} --- duplicates kept --- unless told otherwise, which is
\chref{relational}'s ``a table is not a relation'' showing up in the query
language. Without \sqlkw{DISTINCT}, the second query returns 4{,}000 rows; with
it, one per distinct name and programme pair.

\smallskip
\textbf{And notice the nesting in the third row.} $\sigma$ inside, $\pi$
outside, and SQL writes them in the opposite order from how they run: the
\sqlkw{SELECT} list is applied last, after \sqlkw{WHERE}. Reading SQL in
execution order --- \sqlkw{FROM}, then \sqlkw{WHERE}, then \sqlkw{SELECT} --- is
the single most useful habit this chapter can give you.
\end{examplebox}

\section{The Binary Operators}

\begin{definitionbox}[Product, union, difference, intersection]
The \term{Cartesian product} $R \times S$ pairs every tuple of $R$ with every
tuple of $S$. If $R$ has $m$ tuples and $S$ has $n$, the result has
$m \times n$ tuples and the attributes of both.

\smallskip
\term{Union} $R \cup S$, \term{difference} $R - S$ and \term{intersection}
$R \cap S$ are the set operations of Appendix~B, applied to tuples. All three
require $R$ and $S$ to be \term{union-compatible}: the same number of
attributes, with compatible domains, in the same order.

\smallskip
Intersection is not fundamental: $R \cap S = R - (R - S)$.
\end{definitionbox}

\begin{worked}[How big is a product, and why that matters]
On the seeded database that \texttt{code/schema/seed.py} builds,
\textit{student} has 4{,}000 tuples and \textit{enrolment} has 60{,}037.

\medskip
\textbf{Step 1 --- the size.} $|\mathit{student} \times \mathit{enrolment}|
= 4{,}000 \times 60{,}037 = 240{,}148{,}000$ tuples.

\medskip
\textbf{Step 2 --- how many of them mean anything.} A pair is meaningful only
when the student in it \emph{is} the student the enrolment is about. Each
enrolment has exactly one such student, so the meaningful pairs number exactly
60{,}037.

\medskip
\textbf{Step 3 --- the proportion.} $60{,}037 / 240{,}148{,}000 = 0.025\%$.
One pair in four thousand is true; the other 3{,}999 pair a student with an
enrolment that belongs to somebody else.

\medskip
\textbf{Step 4 --- the fix, in algebra.} Select the meaningful ones:
\[
  \sigma_{\mathit{student.roll\_no} = \mathit{enrolment.roll\_no}}
  (\mathit{student} \times \mathit{enrolment})
\]
That expression \emph{is} a join, and \chref{joins} gives it its own symbol.
Every join in SQL is a product followed by a selection --- in meaning. In
execution it is never computed that way, because 240 million intermediate
tuples to keep 60 thousand is exactly the plan an optimiser exists to avoid,
and \chref{joins} measures how it avoids it.
\end{worked}

\begin{alertbox}[Union-compatibility is checked by position, and that is a trap]
\begin{lstlisting}[language=sql]
SELECT roll_no, name FROM student
UNION
SELECT emp_id,  name FROM instructor;
\end{lstlisting}

This runs. Two attributes each side, both text: union-compatible. The result
is a list of people --- students and staff --- and that is what was meant.

\begin{lstlisting}[language=sql]
SELECT name,    roll_no FROM student
UNION
SELECT emp_id,  name    FROM instructor;
\end{lstlisting}

This \emph{also} runs, and produces a column that holds student names and
employee numbers, and a second column holding roll numbers and instructor
names. SQL checks that the count and types match \emph{by position}; it does
not check that the meanings do, because it cannot.

\smallskip
\textbf{The rule:} in a \sqlkw{UNION}, \sqlkw{EXCEPT} or \sqlkw{INTERSECT},
list the columns in the same order on both sides and give them the same
meaning. The database will hold you to the first half only.
\end{alertbox}

\section{Rename, and a Question Written Three Ways}

\begin{definitionbox}[Rename]
\term{Rename} $\rho_{S}(R)$ gives the relation $R$ the name $S$;
$\rho_{S(B_1,\dots,B_n)}(R)$ also renames its attributes.

\smallskip
It is needed whenever a relation appears twice in one expression --- the
recursive relationship of \chref{cardinality} --- because otherwise
$\mathit{course.course\_code}$ would name two different things. In SQL it is
the \emph{alias}: \texttt{FROM course AS c1, course AS c2}.
\end{definitionbox}

\begin{worked}[``Which students have taken a course but never received an A?'']
\textbf{Step 1 --- in English, the two sets.} Students who have taken
something, minus students who have an A somewhere.

\medskip
\textbf{Step 2 --- each set in algebra.}
\begin{align*}
  T &= \pi_{\mathit{roll\_no}}(\mathit{enrolment}) \\
  A &= \pi_{\mathit{roll\_no}}(\sigma_{\mathit{grade}=\mathrm{'A'}}
        (\mathit{enrolment}))
\end{align*}

\textbf{Step 3 --- combine.} $T - A$. Both are projections onto
\textit{roll\_no}, so they are union-compatible by construction.

\medskip
\textbf{Step 4 --- the SQL.}
\begin{lstlisting}[language=sql]
SELECT roll_no FROM enrolment
EXCEPT
SELECT roll_no FROM enrolment WHERE grade = 'A';
\end{lstlisting}

\medskip
\textbf{Step 5 --- check it against a wrong answer.} The tempting version is
\[
  \pi_{\mathit{roll\_no}}(\sigma_{\mathit{grade} \neq \mathrm{'A'}}
  (\mathit{enrolment}))
\]
--- ``students with a grade that is not A''. That includes every student who
has one B, even if they also have five A's. It answers a different question
and it runs without complaint. On the seeded database the right expression
returns \textbf{551} students; the tempting one returns \textbf{4{,}000} ---
every student in the university, because every one of them has at least one
grade that is not an A.

\smallskip
\textbf{The lesson.} ``Never'' is a difference, not a selection. Questions
with \emph{no}, \emph{never}, \emph{none} or \emph{only} in them almost
always need $-$ or its SQL equivalents, and \chref{subqueries} returns to
exactly this, with \sqlkw{NOT EXISTS} and a measurement.
\end{worked}

\begin{conceptbox}[Algebra is procedural; SQL is declarative --- and why both exist]
An algebra expression says \emph{how}: first select, then project, then take
the difference. Its nesting is an order of operations.

\smallskip
SQL says \emph{what}: the rows satisfying this condition, from these tables.
It deliberately does not say in what order.

\smallskip
The optimiser sits between them. It turns your SQL into one algebra
expression, then uses algebraic laws --- $\sigma$ before $\times$ rather than
after, $\pi$ as early as possible --- to rewrite it into a cheaper one that
computes the same relation. Those laws are true \emph{only because} every
operator is closed and every intermediate result is a relation.

\smallskip
So closure, the first definition in this chapter, is not a curiosity. It is
the reason a database can choose a plan at all, and \chref{indexes} is about
what happens when it chooses badly.
\end{conceptbox}

\begin{pitfall}
\begin{tight}
  \item \textbf{Forgetting \sqlkw{DISTINCT}.} Algebra $\pi$ removes
        duplicates; SQL's \sqlkw{SELECT} list does not.
  \item \textbf{Reading SQL in written order.} It runs \sqlkw{FROM},
        \sqlkw{WHERE}, \sqlkw{SELECT}. Read it that way.
  \item \textbf{Writing ``never'' as a selection.} $\sigma_{\neq}$ answers
        ``at least once something else''. ``Never'' is a difference.
  \item \textbf{Trusting \sqlkw{UNION} to check meaning.} It checks count and
        type by position. Order the columns yourself.
  \item \textbf{Forgetting the size of a product.} $m \times n$. Four thousand
        by sixty thousand is 240 million.
  \item \textbf{Using a relation twice without renaming it.} The attribute
        names become ambiguous. Use an alias.
\end{tight}
\end{pitfall}

\begin{lab}[Compute five algebra expressions, and catch one wrong answer]
\begin{lstlisting}[language=bash]
# --- 1. A populated database, the same rows as the handout ---------
python code/schema/seed.py university.db
sqlite3 university.db
\end{lstlisting}
\begin{lstlisting}[language=sql]
-- --- 2. Selection and projection -------------------------------
.headers on
.mode box
SELECT COUNT(*) FROM student WHERE semester = 2;
SELECT COUNT(*) FROM (SELECT programme FROM student);
SELECT COUNT(*) FROM (SELECT DISTINCT programme FROM student);
-- The last two are both "projection onto programme". Why do they
-- give different numbers, and which one is the algebra's answer?

-- --- 3. The size of a product, without computing it -------------
SELECT (SELECT COUNT(*) FROM student) * (SELECT COUNT(*) FROM enrolment);
-- Do NOT run SELECT COUNT(*) FROM student, enrolment on this
-- database unless you are prepared to wait. Say why.

-- --- 4. A product made meaningful: the selection that is a join -
SELECT COUNT(*) FROM student s, enrolment e
 WHERE s.roll_no = e.roll_no;
-- Compare with the size of enrolment. Explain the match.

-- --- 5. Union, and the positional trap --------------------------
SELECT roll_no AS id, name FROM student
UNION
SELECT emp_id, name FROM instructor
LIMIT 5;
SELECT name, roll_no FROM student
UNION
SELECT emp_id, name FROM instructor
LIMIT 5;
-- Both run. Describe precisely what is wrong with the second.

-- --- 6. "Never an A" -- the right way ---------------------------
SELECT COUNT(*) FROM (
  SELECT roll_no FROM enrolment
  EXCEPT
  SELECT roll_no FROM enrolment WHERE grade = 'A');

-- --- 7. "Never an A" -- the tempting way ------------------------
SELECT COUNT(DISTINCT roll_no) FROM enrolment WHERE grade <> 'A';
-- Write both counts down. One of them is every student in the
-- university. Which, and why? This is the worked example, run.

-- --- 8. Rename: the recursive relationship ----------------------
SELECT c1.title AS course, c2.title AS requires
  FROM prerequisite p
  JOIN course c1 ON c1.course_code = p.course_code
  JOIN course c2 ON c2.course_code = p.requires
 LIMIT 5;
-- course appears twice. Remove the aliases and record what fails.
\end{lstlisting}

\smallskip
\textbf{What to notice.} Steps 6 and 7 differ by one idea --- 551 against
4{,}000 on the seeded data --- and neither one raises an error. Every query you
write from here should be checked against a hand-worked expectation, because
a confident wrong answer looks identical to a confident right one.
\end{lab}

\begin{checkpoint}
\begin{tightnum}
  \item $\pi_{\mathit{programme}}(\mathit{student})$ on 4{,}000 students
        returns 4 tuples, and \texttt{SELECT programme FROM student} returns
        4{,}000 rows. Explain the difference in terms of the definition of a
        relation, and say which keyword reconciles them.
  \item Write, in algebra and in SQL, ``students in semester 2 who are
        enrolled in nothing in FALL2026''. Say which operator expresses
        ``nothing'', and give the tempting wrong version and what it returns.
  \item Every join is a selection over a product. Using the seeded sizes,
        explain why no database computes a join that way, and say which
        property of the algebra lets the optimiser compute something else
        instead and still be correct.
\end{tightnum}
\end{checkpoint}

\begin{chaptersummary}
\begin{tight}
  \item Relational algebra is a set of operators on relations whose every
        result is \textbf{a relation}. That is \textbf{closure}, and it is why
        queries nest.
  \item Six operators are fundamental: \textbf{selection} $\sigma$ (rows),
        \textbf{projection} $\pi$ (columns), \textbf{product} $\times$,
        \textbf{union} $\cup$, \textbf{difference} $-$, \textbf{rename}
        $\rho$. Intersection and every join are built from them.
  \item Algebra projection \textbf{removes duplicates}; SQL's does not unless
        you write \sqlkw{DISTINCT}. A table is a bag; a relation is a set.
  \item SQL runs \sqlkw{FROM}, then \sqlkw{WHERE}, then \sqlkw{SELECT} --- the
        opposite of the order it is written in.
  \item $|R \times S| = |R| \times |S|$. On the seeded data that is 240
        million tuples, of which 0.025\% mean anything. A join is a selection
        over a product \emph{in meaning}, never in execution.
  \item \sqlkw{UNION}, \sqlkw{EXCEPT} and \sqlkw{INTERSECT} check
        union-compatibility \textbf{by position and type}, never by meaning.
  \item ``Never'', ``none'' and ``only'' are \textbf{differences}, not
        selections. The selection version runs and answers a different
        question.
  \item \textbf{Rename} is the SQL alias, and it is required whenever a
        relation appears twice in one query.
  \item Algebra is procedural, SQL declarative; the optimiser turns one into
        the other and rewrites it using laws that hold only because of
        closure.
\end{tight}
\end{chaptersummary}

\begin{reviewq}
  \item \textbf{(MCQ)} Which operator can return fewer tuples than its input
        even when no condition is given? (a)~selection; (b)~projection;
        (c)~rename; (d)~Cartesian product.
  \item \textbf{(MCQ)} $|R| = 30$ and $|S| = 50$. $|R \times S|$ is:
        (a)~80; (b)~30; (c)~1{,}500; (d)~it depends on the keys.
  \item \textbf{(Short)} Define union-compatibility and give two relations
        from the university schema that are union-compatible but whose union
        would be meaningless.
  \item \textbf{(Short)} Show that intersection can be written using only
        difference.
  \item \textbf{(Short)} Why does SQL need \sqlkw{DISTINCT} when relational
        algebra does not?
  \item \textbf{(Applied)} Write in algebra and then in SQL: (a)~the names of
        instructors in the CS department; (b)~course codes that are a
        prerequisite for something but require nothing themselves;
        (c)~students who have received both an A and an F. For each, say which
        operators you used and why.
  \item \textbf{(Applied)} Draw the query tree for
        $\pi_{\mathit{name}}(\sigma_{\mathit{grade}=\mathrm{'A'}}
        (\mathit{student} \times \mathit{enrolment}))$ with the join condition
        added. Then draw an equivalent tree that is cheaper, say which
        algebraic law justifies each change, and estimate from the seeded
        table sizes how many intermediate tuples each tree produces.
\end{reviewq}
